The construction sector is among the world’s largest industries, with annual expenditures of roughly USD 10 trillion on construction goods and services, representing approximately 13\% of global Gross Domestic Product (GDP) in 2023, and is expected to grow to USD 33 trillion by 2040 \citep{Singh2026}. 

Due to their dynamic and demanding context, engineering projects operate under increasingly compressed timelines, complex supply chains, and stringent performance requirements. Within this operational environment, the construction schedule becomes a mission-critical asset, serving as the primary mechanism for resource orchestration, activity coordination, and alignment across multidisciplinary stakeholders.

Scheduling problems in construction projects lead to negative outcomes, including delays and rework \citep{Kline2025}. Despite its critical importance in construction projects, scheduling is often underemphasized in formal engineering education. Consequently, construction and project management professionals tend to acquire skills in developing workable schedules primarily through on-the-job experience \citep{Vahidi2025}.

In this paper, we develop a digital twin framework to assess risk in construction scheduling. We explore the traditional PERT (Program Evaluation and Review Technique) and CPM (Critical Path Method) methods within a probabilistic framework to estimate scheduling risk metrics. Uncertainty and dependence of the duration of each project's activities are taken into account using Monte Carlo simulation and Bayesian Networks. 

The paper is structured as follows. In the next section, we present the key tools and techniques for assessing scheduling risks. We then briefly discuss the digital twin framework and contextualize the proposed model. We describe the architecture of the proposed framework and discuss the results obtained from its application to assess scheduling risks, using a toy problem in civil engineering. We present characteristics of the implemented system, including details of the simulation mechanism and risk analysis. Finally, we make concluding remarks, highlighting the applications and limitations of the study.

